% Generated by SIGMA; compile with Tectonic or another TeX engine.
\documentclass[11pt]{article}
\usepackage[T1]{fontenc}
\usepackage[utf8]{inputenc}
\usepackage[margin=1in]{geometry}
\usepackage{amsmath,amssymb,booktabs,longtable}
\usepackage{array}
\usepackage{enumitem}
\setlist{nosep,leftmargin=*}
\usepackage{microtype}
\usepackage{titlesec}
\usepackage{xcolor}
\definecolor{blue}{HTML}{24485E}
\titleformat{\section}{\Large\bfseries\sffamily\color{blue}}{\thesection}{0.8em}{}
\titleformat{\subsection}{\large\bfseries\sffamily}{\thesubsection}{0.7em}{}
\titleformat{\subsubsection}{\normalsize\bfseries\sffamily\color{blue}}{\thesubsubsection}{0.6em}{}
\setlength{\parskip}{0.55em}
\setlength{\parindent}{0pt}
\title{Nonconsecutive Subset Sums Modulo 3 v1}
\author{SIGMA Research Record}
\date{\today}
\begin{document}
\maketitle
\begin{quote}\small In the future we plan on using AI to produce a typical academic paper based on the results of the research.\end{quote}
\section{Project}
\textbf{Interpretation:} Agent-generated findings are provisional unless explicitly verified; reviewer agreement is not proof.\par
\textbf{Research problem:} For each integer n \textgreater{}= 0, let a\_n count the subsets S of \{1,...,n\} with no two consecutive elements such that the sum of the elements of S is divisible by 3. The empty subset counts. Derive and rigorously prove a constant-coefficient linear recurrence with integer coefficients for a\_n, giving sufficient initial values. Determine a\_12 exactly. Finally prove that a\_n/F\_\{n+2\} tends to 1/3, where F\_0=0 and F\_1=1 are the Fibonacci numbers. Computational pattern matching alone does not count as a proof.\par
\textbf{Project ID:} 1e8625bf-12fd-4148-8db9-8878895466fd\par
\textbf{Research type:} Mathematics\par
\section{Strategies}
\subsection{Periodic transfer matrix (saved)}
Use a dynamic-programming state (r,b), with r in \{0,1,2\} the sum of selected indices modulo 3 and b indicating whether the most recently processed index was selected. On processing i, either omit i, which resets b=0, or include i only if b=0, updating r to r+i mod 3. Encode this as a 6-by-6 transition matrix T\_i. Because T\_\{i+3\}=T\_i, calculate the three-step matrix and seek a polynomial identity satisfied by each residue-class subsequence of a\_n. First tests: independently enumerate admissible subsets for n=0,...,8 and compare all three residues. Proof obligations: derive an exact constant-coefficient recurrence, list sufficient initial values, and bound nontrivial residue modes strongly enough to prove asymptotic equidistribution. If the matrix characteristic polynomial is unnecessarily large, look for invariant subspaces or symmetry reduction; do not infer a recurrence solely from samples.\par
\subsection{Roots-of-unity weighted Fibonacci polynomials (saved)}
Define P\_n(q)=sum q\textasciicircum{}\{sum\_\{i in S\} i\} over nonconsecutive subsets S of \{1,...,n\}. Separating subsets according to whether n is present yields P\_n(q)=P\_\{n-1\}(q)+q\textasciicircum{}n P\_\{n-2\}(q), with P\_0(q)=1 and P\_1(q)=1+q. For a primitive cube root omega, use a\_n=(P\_n(1)+P\_n(omega)+P\_n(omega\textasciicircum{}2))/3. At q=1 use ordinary Fibonacci growth; at omega and omega\textasciicircum{}2 examine the periodic 2-by-2 transfer matrices to derive a recurrence over n -\textgreater{} n+3. Combine their annihilating polynomials into a proved integer-coefficient recurrence for a\_n. First tests: evaluate P\_n modulo q\textasciicircum{}3-1 for n\textless{}=12 and cross-check the six-state DP. Proof obligation: control both complex terms relative to P\_n(1) for the 1/3 limit. Do not confuse pointwise evaluation at a root of unity with extracting a coefficient.\par
\subsection{Gap removal and Gaussian binomial sums (saved)}
For a nonconsecutive k-subset 1\textless{}=i\_1\textless{}...\textless{}i\_k\textless{}=n, set y\_j=i\_j-(j-1). Then 1\textless{}=y\_1\textless{}...\textless{}y\_k\textless{}=n-k+1 and sum i\_j=sum y\_j+k(k-1)/2. The ordinary k-subset weight enumerator gives P\_n(q)=sum\_\{k\textgreater{}=0\} q\textasciicircum{}\{k\textasciicircum{}2\} [n-k+1 choose k]\_q, with Gaussian q-binomials and k\textless{}=ceil(n/2). Derive this carefully, verify low-n polynomials, then evaluate its 3-dissection using q=1,omega,omega\textasciicircum{}2; explore q-Lucas identities at cube roots or elementary residue-block arguments. Seek an independent demonstration of the recurrence and asymptotic density, or a clear algebraic reduction to one of those statements. Discriminating test: check any proposed closed form against the weighted-Fibonacci recurrence for several complete periods, not just for n divisible by 3.\par
\subsection{Close the phase-matrix proof obligation (saved)}
Use Session 1's n=2 counterexample to exclude the short guessed recurrence. Take Session 2's order-12 relation as a candidate only. Multiply each ordered three-phase matrix symbolically, handle the n=0 boundary, and bound nontrivial-mode growth before asserting the 1/3 limit. Preserve Session 5's corrected q-binomial normalization. Evaluator-authored continuation; no new research run.\par
\section{Research sessions}
\subsection{Session \#1}
\begin{description}[style=nextline,font=\normalfont\bfseries,leftmargin=0pt,labelsep=0pt]
\item[Execution status] ended
\item[Outcome] Mixed · partial progress, unresolved
\item[Strategy] Periodic transfer matrix
\item[Started] 2026-10-23T21:00:13.199709+00:00
\end{description}
\textbf{Outcome summary:} The tentative recurrence a\_(n+3)=a\_(n+2)+a\_n is false; the original problem remains open.\par
\subsection{Conclusion}
\#\#\# Historical session

No Conclusion was saved for this session. Its existing outcome summary and detailed evidence remain available below.

Structured argument steps were not recorded; saved narrative is available for inspection in Reasoning / proof and the research artifacts. Its completeness and correctness remain unassessed.\par
Outcome caveats:\par
\begin{itemize}
\item Agent-reported evidence is unverified; agreement is not independent verification.
\item The saved argument remains incomplete; no solution is claimed.
\item No general recurrence or asymptotic proof has been established.
\item Open obligations remain; inspect the saved artifacts before relying on this result.
\end{itemize}
\textbf{Verification:} Incomplete\par
Structured argument steps were not recorded; saved narrative is available for inspection in Reasoning / proof and the research artifacts. Its completeness and correctness remain unassessed.\par
\subsection{Outcome metrics}
\subsection{Operational metrics}
\textbf{Participating agents:} Investigator, Skeptic\par
\textbf{Model calls by status:} 2 dispatched · 2 completed · 0 rejected · 0 timed out · 0 interrupted / cancelled · 0 failed · 0 in progress · 0 not dispatched · 0 status unavailable\par
Tokens · input / output / cached / total: — (0/2 calls report token usage) / — (0/2 calls report token usage) / — (0/2 calls include cached-token metadata) / — (token usage incomplete for 2 call(s))\par
Wall-clock / agent call time: 0s wall · Investigator 0s across 1 call(s) · Skeptic 0s across 1 call(s)\par
\textbf{Python runs by status:} 1 completed\par
\textbf{SageMath runs by status:} 0 executions\par
\textbf{Artifacts and storage:} 1 finding(s) · 0 Python file(s) (0 B) · session record 29,539 B\par
Provider failures / retries: 0 provider failure(s) (0 rejected, 0 timed out, 0 other) · 0 interrupted / cancelled · 0 later attempt(s); no automatic retries\par
\textbf{Estimated monetary cost:} — (complete usage metadata or applicable price unavailable)\par
\subsection{Research-quality metrics}
\textbf{Strategies investigated:} Periodic transfer matrix\par
\textbf{Research cycles:} 1 completed / 1 started\par
\textbf{Outcome classification:} Mixed · partial progress, unresolved\par
\textbf{Verification status:} Incomplete\par
\textbf{Open proof obligations:} 5 · Derive a general recurrence from the periodic transition matrices.; Prove the requested asymptotic limit. (+3 more)\par
\textbf{Evidence-backed conclusions:} 1 conclusion(s) · 1 cited finding(s)\par
\textbf{Strategy contribution:} not used: 2 · productive: 1\par
\textbf{Unresolved objections:} 1 · No general recurrence or asymptotic proof has been established.\par
\textbf{Claim scope:} not recorded\par
\textbf{Solution claim:} The tentative recurrence a\_(n+3)=a\_(n+2)+a\_n is false; the original problem remains open.\par
\textbf{Argument:} A small exact enumeration rejects one candidate rule at n=2. The counterexample is informative but establishes no recurrence for a\_n.\par
\textbf{Skeptic response:} Deterministic fixture result (not live-model output). This Skeptic response is a deterministic local fixture, not a live-model review. The computed values are [1,1,1,2,3,5,8,12,19]. The candidate relation holds for n=0,1 but fails at n=2: a\_5=5 whereas a\_4+a\_2=3+1=4. This refutes only the tentative recurrence, not the original counting problem.\par
\textbf{Assumptions reviewed:} The enumeration includes the empty subset and tests adjacency on the bit mask.\par
\textbf{Proof obligations considered:} Derive a general recurrence from the periodic transition matrices.; Prove the requested asymptotic limit.\par
\textbf{Still open:} Check the proposed relation beyond the first two transitions.; Derive a general transition identity before claiming any recurrence.; Replace the false pattern with a general transfer-matrix identity.\par
\textbf{Solution evidence references:} 1cc6ec50-87c0-46f4-b1dd-a82857d74f9e\par
\textbf{Unresolved objections:} No general recurrence or asymptotic proof has been established.\par
\textbf{Resolved review notes:} None recorded\par
\textbf{Skeptic objection:} The proposed three-step recurrence has the explicit n=2 counterexample a\_5=5 != a\_4+a\_2=4.\par
\textbf{Saved verification code:} incomplete\par
\textbf{Strategy finding:} Periodic transfer matrix · partial\par
\textbf{Strategy execution status:} ended\par
\textbf{Strategy failure detail:} Stop budget exhausted\par
\textbf{Claim:} The tentative recurrence a\_(n+3)=a\_(n+2)+a\_n is false; the original problem remains open.\par
\textbf{Source artifacts:} 1cc6ec50-87c0-46f4-b1dd-a82857d74f9e\par
\textbf{Objection:} The proposed three-step recurrence has the explicit n=2 counterexample a\_5=5 != a\_4+a\_2=4.\par
\textbf{Open obligation:} Check the proposed relation beyond the first two transitions.\par
\textbf{Open obligation:} Derive a general transition identity before claiming any recurrence.\par
\textbf{Open obligation:} Replace the false pattern with a general transfer-matrix identity.\par
\subsubsection*{Strategy assessments}
\begin{longtable}{>{\raggedright\arraybackslash}p{.16\linewidth}>{\raggedright\arraybackslash}p{.15\linewidth}>{\raggedright\arraybackslash}p{.40\linewidth}>{\raggedright\arraybackslash}p{.17\linewidth}}
\toprule Strategy & Assessment & Rationale & Evidence \\
\midrule\endhead
Periodic transfer matrix & Productive & The falsified rule narrows the investigation and motivates an algebraic transition-matrix derivation. & 1cc6ec50-87c0-46f4-b1dd-a82857d74f9e \\
Roots-of-unity weighted Fibonacci polynomials & Not used & No session artifact is linked to this strategy. & No artifact references \\
Gap removal and Gaussian binomial sums & Not used & No session artifact is linked to this strategy. & No artifact references \\
\bottomrule
\end{longtable}
\paragraph{Artifact 1 · Unverified research note. Any Investigator Python or SageMath execution, including its source, status, output, and captured artifacts, is recorded below and is not proof by itself.}
\textbf{Claim:} A tentative three-step rule a\_(n+3)=a\_(n+2)+a\_n fits the first two transitions, but has not been justified.\par
\textbf{Method:} Start with the six states (sum mod 3, previous-selected bit), enumerate small n independently, and inspect the three-step count table. This scripted fixture deliberately tests a plausible early pattern; it is not live-model inference.\par
\textbf{Assumptions:} Indices are one-based in subset sums.; The empty subset contributes at n=0.\par
\textbf{Evidence status:} Deterministic local fixture response for UI review; not live-model inference. Python enumeration is bounded to n\textless{}=8 and can only test this candidate.\par
\textbf{Skeptic response:} Deterministic fixture result (not live-model output). This Skeptic response is a deterministic local fixture, not a live-model review. The computed values are [1,1,1,2,3,5,8,12,19]. The candidate relation holds for n=0,1 but fails at n=2: a\_5=5 whereas a\_4+a\_2=3+1=4. This refutes only the tentative recurrence, not the original counting problem.\par
\textbf{Objections:} The proposed three-step recurrence has the explicit n=2 counterexample a\_5=5 != a\_4+a\_2=4.\par
\textbf{Decision:} stop\_budget\_exhausted\par
\textbf{Open obligations:} Check the proposed relation beyond the first two transitions.; Derive a general transition identity before claiming any recurrence.; Replace the false pattern with a general transfer-matrix identity.\par
\textbf{Next question:} Not recorded\par
\textbf{Terminal reason:} stop\_budget\_exhausted\par
\subsection{Session \#2}
\begin{description}[style=nextline,font=\normalfont\bfseries,leftmargin=0pt,labelsep=0pt]
\item[Execution status] ended
\item[Outcome] Mixed · partial progress, unresolved
\item[Strategy] Roots-of-unity weighted Fibonacci polynomials
\item[Started] 2026-10-23T21:00:13.734982+00:00
\end{description}
\textbf{Outcome summary:} Deterministic fixture result (not live-model output): the candidate step-three recurrence and a\_12=128 are consistent with finite computation, but no global proof is recorded.\par
\subsection{Conclusion}
\#\#\# Historical session

No Conclusion was saved for this session. Its existing outcome summary and detailed evidence remain available below.

Structured argument steps were not recorded; saved narrative is available for inspection in Reasoning / proof and the research artifacts. Its completeness and correctness remain unassessed.\par
Outcome caveats:\par
\begin{itemize}
\item Agent-reported evidence is unverified; agreement is not independent verification.
\item The saved argument remains incomplete; no solution is claimed.
\item General matrix identity and asymptotic bound are not independently verified.
\item Open obligations remain; inspect the saved artifacts before relying on this result.
\end{itemize}
\textbf{Verification:} Incomplete\par
Structured argument steps were not recorded; saved narrative is available for inspection in Reasoning / proof and the research artifacts. Its completeness and correctness remain unassessed.\par
\subsection{Outcome metrics}
\subsection{Operational metrics}
\textbf{Participating agents:} Investigator, Skeptic\par
\textbf{Model calls by status:} 2 dispatched · 2 completed · 0 rejected · 0 timed out · 0 interrupted / cancelled · 0 failed · 0 in progress · 0 not dispatched · 0 status unavailable\par
Tokens · input / output / cached / total: — (0/2 calls report token usage) / — (0/2 calls report token usage) / — (0/2 calls include cached-token metadata) / — (token usage incomplete for 2 call(s))\par
Wall-clock / agent call time: 0s wall · Investigator 0s across 1 call(s) · Skeptic 0s across 1 call(s)\par
\textbf{Python runs by status:} 1 completed\par
\textbf{SageMath runs by status:} 0 executions\par
\textbf{Artifacts and storage:} 1 finding(s) · 0 Python file(s) (0 B) · session record 37,860 B\par
Provider failures / retries: 0 provider failure(s) (0 rejected, 0 timed out, 0 other) · 0 interrupted / cancelled · 0 later attempt(s); no automatic retries\par
\textbf{Estimated monetary cost:} — (complete usage metadata or applicable price unavailable)\par
\subsection{Research-quality metrics}
\textbf{Strategies investigated:} Roots-of-unity weighted Fibonacci polynomials\par
\textbf{Research cycles:} 1 completed / 1 started\par
\textbf{Outcome classification:} Mixed · partial progress, unresolved\par
\textbf{Verification status:} Incomplete\par
\textbf{Open proof obligations:} 7 · Multiply the three phase matrices and verify trace 1, determinant -1.; Combine characteristic polynomials with the q=1 Fibonacci component. (+5 more)\par
\textbf{Evidence-backed conclusions:} 1 conclusion(s) · 1 cited finding(s)\par
\textbf{Strategy contribution:} not used: 2 · productive: 1\par
\textbf{Unresolved objections:} 1 · General matrix identity and asymptotic bound are not independently verified.\par
\textbf{Claim scope:} not recorded\par
\textbf{Solution claim:} Deterministic fixture result (not live-model output): the candidate step-three recurrence and a\_12=128 are consistent with finite computation, but no global proof is recorded.\par
\textbf{Argument:} The root-of-unity decomposition gives a plausible order-12 recurrence. The saved computation is corroboration only; boundary and asymptotic proof obligations remain.\par
\textbf{Skeptic response:} This Skeptic response is a deterministic local fixture, not a live-model review. The finite DP supports a\_12=128 and all tested recurrence residuals vanish, but neither computation proves a general recurrence. The written matrix argument must show the product order and n=0 boundary explicitly. The claimed limit additionally needs a bound on the nontrivial modes; observed ratios are not a proof.\par
\textbf{Assumptions reviewed:} F\_(n+2) is the unweighted count of nonconsecutive subsets.; The two nontrivial root-of-unity values are complex conjugates.\par
\textbf{Proof obligations considered:} Multiply the three phase matrices and verify trace 1, determinant -1.; Combine characteristic polynomials with the q=1 Fibonacci component.; Establish the asymptotic ratio rigorously.\par
\textbf{Still open:} Show the three-step matrix product for each residue class, including the n=0 boundary.; Prove the omega terms grow strictly slower than F\_(n+2), rather than relying on finite ratios.; Verify the omega matrix identity at n=0 and for all later residue blocks.; Bound the omega and conjugate contributions relative to Fibonacci growth.\par
\textbf{Solution evidence references:} ccc1045f-479d-4018-b8ed-42a60b94e567\par
\textbf{Unresolved objections:} General matrix identity and asymptotic bound are not independently verified.\par
\textbf{Resolved review notes:} None recorded\par
\textbf{Skeptic objection:} A finite range cannot establish the recurrence for all n.\par
\textbf{Skeptic objection:} The n=0 transition is outside the stated n\textgreater{}=2 weighted recurrence and needs a separate check.\par
\textbf{Skeptic objection:} No analytic estimate yet proves the nontrivial terms are negligible compared with F\_(n+2).\par
\textbf{Saved verification code:} incomplete\par
\textbf{Strategy finding:} Roots-of-unity weighted Fibonacci polynomials · partial\par
\textbf{Strategy execution status:} ended\par
\textbf{Strategy failure detail:} Stop budget exhausted\par
\textbf{Claim:} Deterministic fixture result (not live-model output): the candidate step-three recurrence and a\_12=128 are consistent with finite computation, but no global proof is recorded.\par
\textbf{Source artifacts:} ccc1045f-479d-4018-b8ed-42a60b94e567\par
\textbf{Objection:} A finite range cannot establish the recurrence for all n.\par
\textbf{Objection:} The n=0 transition is outside the stated n\textgreater{}=2 weighted recurrence and needs a separate check.\par
\textbf{Objection:} No analytic estimate yet proves the nontrivial terms are negligible compared with F\_(n+2).\par
\textbf{Open obligation:} Show the three-step matrix product for each residue class, including the n=0 boundary.\par
\textbf{Open obligation:} Prove the omega terms grow strictly slower than F\_(n+2), rather than relying on finite ratios.\par
\textbf{Open obligation:} Verify the omega matrix identity at n=0 and for all later residue blocks.\par
\textbf{Open obligation:} Bound the omega and conjugate contributions relative to Fibonacci growth.\par
\subsubsection*{Strategy assessments}
\begin{longtable}{>{\raggedright\arraybackslash}p{.16\linewidth}>{\raggedright\arraybackslash}p{.15\linewidth}>{\raggedright\arraybackslash}p{.40\linewidth}>{\raggedright\arraybackslash}p{.17\linewidth}}
\toprule Strategy & Assessment & Rationale & Evidence \\
\midrule\endhead
Periodic transfer matrix & Not used & No session artifact is linked to this strategy. & No artifact references \\
Roots-of-unity weighted Fibonacci polynomials & Productive & It supplies an exact candidate recurrence and a\_12 while keeping finite computation separate from proof. & ccc1045f-479d-4018-b8ed-42a60b94e567 \\
Gap removal and Gaussian binomial sums & Not used & No session artifact is linked to this strategy. & No artifact references \\
\bottomrule
\end{longtable}
\paragraph{Artifact 1 · Unverified research note. Any Investigator Python or SageMath execution, including its source, status, output, and captured artifacts, is recorded below and is not proof by itself.}
\textbf{Claim:} Deterministic fixture result (not live-model output): the root-of-unity filter suggests a\_(n+12)-5a\_(n+9)+2a\_(n+6)+5a\_(n+3)+a\_n=0 and gives a\_12=128.\par
\textbf{Method:} Write P\_n(q)=P\_(n-1)(q)+q\textasciicircum{}n P\_(n-2)(q), with P\_n(1)=F\_(n+2), and extract the coefficient of residue 0 by averaging P\_n(1), P\_n(omega), and P\_n(omega\textasciicircum{}2). Over each three-step block the omega transfer matrix has trace 1 and determinant -1, hence its subsequences satisfy x\textasciicircum{}2-x-1; the ordinary Fibonacci component advances by x\textasciicircum{}2-4x-1. Their product yields the stated recurrence. This derivation is a scripted fixture narrative and remains unverified in SIGMA.\par
\textbf{Assumptions:} omega is a primitive cube root of unity.; The weighted recurrence is applied with its n=0,1 initial values handled separately.\par
\textbf{Evidence status:} Deterministic local fixture response for UI review; not live-model inference. Python DP checks terms and recurrence residuals only over a finite range; it does not prove the matrix identity or asymptotic bound.\par
\textbf{Skeptic response:} This Skeptic response is a deterministic local fixture, not a live-model review. The finite DP supports a\_12=128 and all tested recurrence residuals vanish, but neither computation proves a general recurrence. The written matrix argument must show the product order and n=0 boundary explicitly. The claimed limit additionally needs a bound on the nontrivial modes; observed ratios are not a proof.\par
\textbf{Objections:} A finite range cannot establish the recurrence for all n.; The n=0 transition is outside the stated n\textgreater{}=2 weighted recurrence and needs a separate check.; No analytic estimate yet proves the nontrivial terms are negligible compared with F\_(n+2).\par
\textbf{Decision:} stop\_budget\_exhausted\par
\textbf{Open obligations:} Show the three-step matrix product for each residue class, including the n=0 boundary.; Prove the omega terms grow strictly slower than F\_(n+2), rather than relying on finite ratios.; Verify the omega matrix identity at n=0 and for all later residue blocks.; Bound the omega and conjugate contributions relative to Fibonacci growth.\par
\textbf{Next question:} Not recorded\par
\textbf{Terminal reason:} stop\_budget\_exhausted\par
\subsection{Session \#3}
\begin{description}[style=nextline,font=\normalfont\bfseries,leftmargin=0pt,labelsep=0pt]
\item[Execution status] ended
\item[Outcome] Mixed · partial progress, unresolved
\item[Strategy] Gap removal and Gaussian binomial sums
\item[Started] 2026-10-23T21:00:14.267271+00:00
\end{description}
\textbf{Outcome summary:} Deterministic fixture result (not live-model output): the gap-removal bijection gives P\_n(q)=sum\_k q\textasciicircum{}(k\textasciicircum{}2) [n-k+1 choose k]\_q, with finite coefficient agreement through n=18.\par
\subsection{Conclusion}
\#\#\# Historical session

No Conclusion was saved for this session. Its existing outcome summary and detailed evidence remain available below.

Structured argument steps were not recorded; saved narrative is available for inspection in Reasoning / proof and the research artifacts. Its completeness and correctness remain unassessed.\par
Outcome caveats:\par
\begin{itemize}
\item Agent-reported evidence is unverified; agreement is not independent verification.
\item The saved argument remains incomplete; no solution is claimed.
\item No complete proof of the recurrence or the limit is recorded.
\item Open obligations remain; inspect the saved artifacts before relying on this result.
\end{itemize}
\textbf{Verification:} Incomplete\par
Structured argument steps were not recorded; saved narrative is available for inspection in Reasoning / proof and the research artifacts. Its completeness and correctness remain unassessed.\par
\subsection{Outcome metrics}
\subsection{Operational metrics}
\textbf{Participating agents:} Investigator, Skeptic\par
\textbf{Model calls by status:} 2 dispatched · 2 completed · 0 rejected · 0 timed out · 0 interrupted / cancelled · 0 failed · 0 in progress · 0 not dispatched · 0 status unavailable\par
Tokens · input / output / cached / total: — (0/2 calls report token usage) / — (0/2 calls report token usage) / — (0/2 calls include cached-token metadata) / — (token usage incomplete for 2 call(s))\par
Wall-clock / agent call time: 0s wall · Investigator 0s across 1 call(s) · Skeptic 0s across 1 call(s)\par
\textbf{Python runs by status:} 1 completed\par
\textbf{SageMath runs by status:} 0 executions\par
\textbf{Artifacts and storage:} 1 finding(s) · 0 Python file(s) (0 B) · session record 42,793 B\par
Provider failures / retries: 0 provider failure(s) (0 rejected, 0 timed out, 0 other) · 0 interrupted / cancelled · 0 later attempt(s); no automatic retries\par
\textbf{Estimated monetary cost:} — (complete usage metadata or applicable price unavailable)\par
\subsection{Research-quality metrics}
\textbf{Strategies investigated:} Gap removal and Gaussian binomial sums\par
\textbf{Research cycles:} 1 completed / 1 started\par
\textbf{Outcome classification:} Mixed · partial progress, unresolved\par
\textbf{Verification status:} Incomplete\par
\textbf{Open proof obligations:} 6 · Derive a general linear recurrence with integer coefficients.; Prove the asymptotic density limit rather than extrapolating from finite coefficients. (+4 more)\par
\textbf{Evidence-backed conclusions:} 1 conclusion(s) · 1 cited finding(s)\par
\textbf{Strategy contribution:} not used: 2 · productive: 1\par
\textbf{Unresolved objections:} 1 · No complete proof of the recurrence or the limit is recorded.\par
\textbf{Claim scope:} not recorded\par
\textbf{Solution claim:} Deterministic fixture result (not live-model output): the gap-removal bijection gives P\_n(q)=sum\_k q\textasciicircum{}(k\textasciicircum{}2) [n-k+1 choose k]\_q, with finite coefficient agreement through n=18.\par
\textbf{Argument:} The substitution y\_j=i\_j-(j-1) is an exact bijection to ordinary k-subsets and explains the q\textasciicircum{}(k\textasciicircum{}2) weight. Its q-binomial form yields a second representation and a\_12=128, but the requested general recurrence and limiting ratio remain unresolved in this record.\par
\textbf{Skeptic response:} This Skeptic response is a deterministic local fixture, not a live-model review. The gap-removal bijection and q-binomial formula are correct as a combinatorial identity, and exact polynomial arithmetic reproduces the first 19 polynomials. This is a useful independent encoding, but the finite check does not finish a recurrence proof or the asymptotic limit.\par
\textbf{Assumptions reviewed:} The inverse map i\_j=y\_j+(j-1) preserves the nonconsecutive condition.; Gaussian q-binomial coefficients use integer polynomial coefficients.\par
\textbf{Proof obligations considered:} Derive a general linear recurrence with integer coefficients.; Prove the asymptotic density limit rather than extrapolating from finite coefficients.\par
\textbf{Still open:} Use a correctly stated q-Lucas or residue-class argument to obtain a general recurrence from the identity.; Supply a rigorous asymptotic estimate for the nontrivial cube-root terms.; Derive the recurrence globally from a justified root-of-unity evaluation.; Prove a\_n/F\_(n+2) tends to 1/3.\par
\textbf{Solution evidence references:} 02ca743b-d66f-41c0-8692-34579ea5adfa\par
\textbf{Unresolved objections:} No complete proof of the recurrence or the limit is recorded.\par
\textbf{Resolved review notes:} None recorded\par
\textbf{Saved verification code:} incomplete\par
\textbf{Strategy finding:} Gap removal and Gaussian binomial sums · partial\par
\textbf{Strategy execution status:} ended\par
\textbf{Strategy failure detail:} Stop productive conclusion\par
\textbf{Claim:} Deterministic fixture result (not live-model output): the gap-removal bijection gives P\_n(q)=sum\_k q\textasciicircum{}(k\textasciicircum{}2) [n-k+1 choose k]\_q, with finite coefficient agreement through n=18.\par
\textbf{Source artifacts:} 02ca743b-d66f-41c0-8692-34579ea5adfa\par
\textbf{Open obligation:} Use a correctly stated q-Lucas or residue-class argument to obtain a general recurrence from the identity.\par
\textbf{Open obligation:} Supply a rigorous asymptotic estimate for the nontrivial cube-root terms.\par
\textbf{Open obligation:} Derive the recurrence globally from a justified root-of-unity evaluation.\par
\textbf{Open obligation:} Prove a\_n/F\_(n+2) tends to 1/3.\par
\subsubsection*{Strategy assessments}
\begin{longtable}{>{\raggedright\arraybackslash}p{.16\linewidth}>{\raggedright\arraybackslash}p{.15\linewidth}>{\raggedright\arraybackslash}p{.40\linewidth}>{\raggedright\arraybackslash}p{.17\linewidth}}
\toprule Strategy & Assessment & Rationale & Evidence \\
\midrule\endhead
Periodic transfer matrix & Not used & No session artifact is linked to this strategy. & No artifact references \\
Roots-of-unity weighted Fibonacci polynomials & Not used & No session artifact is linked to this strategy. & No artifact references \\
Gap removal and Gaussian binomial sums & Productive & The bijection is a genuinely different exact encoding and provides useful finite cross-checks without overstating them as a full solution. & 02ca743b-d66f-41c0-8692-34579ea5adfa \\
\bottomrule
\end{longtable}
\paragraph{Artifact 1 · Unverified research note. Any Investigator Python or SageMath execution, including its source, status, output, and captured artifacts, is recorded below and is not proof by itself.}
\textbf{Claim:} The gap-removal bijection gives P\_n(q)=sum\_k q\textasciicircum{}(k\textasciicircum{}2) [n-k+1 choose k]\_q; exact polynomial checks agree with the weighted recurrence through n=18.\par
\textbf{Method:} For a nonconsecutive k-subset i\_1\textless{}...\textless{}i\_k, set y\_j=i\_j-(j-1). Then y is an ordinary k-subset of \{1,...,n-k+1\} and sum(i\_j)=sum(y\_j)+k(k-1)/2. The Gaussian-binomial generating function contributes q\textasciicircum{}(k(k-1)/2), giving the q\textasciicircum{}(k\textasciicircum{}2) factor. This scripted fixture narrative is not live-model inference.\par
\textbf{Assumptions:} The shifted positions remain strictly increasing and range from 1 to n-k+1.; Polynomial equality is checked with integer coefficients before specializing q.\par
\textbf{Evidence status:} Deterministic local fixture response for UI review; not live-model inference. Exact Python polynomial arithmetic checks only n\textless{}=18 and does not prove a density limit.\par
\textbf{Skeptic response:} This Skeptic response is a deterministic local fixture, not a live-model review. The gap-removal bijection and q-binomial formula are correct as a combinatorial identity, and exact polynomial arithmetic reproduces the first 19 polynomials. This is a useful independent encoding, but the finite check does not finish a recurrence proof or the asymptotic limit.\par
\textbf{Objections:} None recorded\par
\textbf{Decision:} stop\_productive\_conclusion\par
\textbf{Open obligations:} Use a correctly stated q-Lucas or residue-class argument to obtain a general recurrence from the identity.; Supply a rigorous asymptotic estimate for the nontrivial cube-root terms.; Derive the recurrence globally from a justified root-of-unity evaluation.; Prove a\_n/F\_(n+2) tends to 1/3.\par
\textbf{Next question:} Not recorded\par
\textbf{Terminal reason:} stop\_productive\_conclusion\par
\subsection{Session \#4}
\begin{description}[style=nextline,font=\normalfont\bfseries,leftmargin=0pt,labelsep=0pt]
\item[Execution status] ended
\item[Outcome] Mixed · partial progress, unresolved
\item[Strategy] Periodic transfer matrix
\item[Started] 2026-10-23T21:27:59.590672+00:00
\end{description}
\textbf{Outcome summary:} Deterministic fixture result (not live-model output): a finite exact DP through n=60 is consistent with the candidate order-12 step-three recurrence; the general recurrence and limiting ratio remain unproved in this record.\par
\subsection{Conclusion}
\#\#\# Historical session

No Conclusion was saved for this session. Its existing outcome summary and detailed evidence remain available below.

Structured argument steps were not recorded; saved narrative is available for inspection in Reasoning / proof and the research artifacts. Its completeness and correctness remain unassessed.\par
Outcome caveats:\par
\begin{itemize}
\item Agent-reported evidence is unverified; agreement is not independent verification.
\item The saved argument remains incomplete; no solution is claimed.
\item The general recurrence has not been derived from the matrices in this artifact.
\item The requested limiting density remains unproved.
\item Open obligations remain; inspect the saved artifacts before relying on this result.
\end{itemize}
\textbf{Verification:} Incomplete\par
Structured argument steps were not recorded; saved narrative is available for inspection in Reasoning / proof and the research artifacts. Its completeness and correctness remain unassessed.\par
\subsection{Outcome metrics}
\subsection{Operational metrics}
\textbf{Participating agents:} Investigator, Skeptic\par
\textbf{Model calls by status:} 2 dispatched · 2 completed · 0 rejected · 0 timed out · 0 interrupted / cancelled · 0 failed · 0 in progress · 0 not dispatched · 0 status unavailable\par
Tokens · input / output / cached / total: — (0/2 calls report token usage) / — (0/2 calls report token usage) / — (0/2 calls include cached-token metadata) / — (token usage incomplete for 2 call(s))\par
Wall-clock / agent call time: 0s wall · Investigator 0s across 1 call(s) · Skeptic 0s across 1 call(s)\par
\textbf{Python runs by status:} 1 completed\par
\textbf{SageMath runs by status:} 0 executions\par
\textbf{Artifacts and storage:} 1 finding(s) · 0 Python file(s) (0 B) · session record 51,697 B\par
Provider failures / retries: 0 provider failure(s) (0 rejected, 0 timed out, 0 other) · 0 interrupted / cancelled · 0 later attempt(s); no automatic retries\par
\textbf{Estimated monetary cost:} — (complete usage metadata or applicable price unavailable)\par
\subsection{Research-quality metrics}
\textbf{Strategies investigated:} Periodic transfer matrix\par
\textbf{Research cycles:} 1 completed / 1 started\par
\textbf{Outcome classification:} Mixed · partial progress, unresolved\par
\textbf{Verification status:} Incomplete\par
\textbf{Open proof obligations:} 8 · Show the ordered transfer-matrix polynomial identity symbolically.; Establish a rigorous asymptotic bound for the nontrivial modes. (+6 more)\par
\textbf{Evidence-backed conclusions:} 1 conclusion(s) · 1 cited finding(s)\par
\textbf{Strategy contribution:} not used: 2 · productive: 1\par
\textbf{Unresolved objections:} 2 · The general recurrence has not been derived from the matrices in this artifact.; The requested limiting density remains unproved.\par
\textbf{Claim scope:} not recorded\par
\textbf{Solution claim:} Deterministic fixture result (not live-model output): a finite exact DP through n=60 is consistent with the candidate order-12 step-three recurrence; the general recurrence and limiting ratio remain unproved in this record.\par
\textbf{Argument:} The revisit uses the counterexample to avoid repeating the false short recurrence and uses the other sessions as candidate/independent-comparison sources. It adds bounded exact computational evidence but does not close any global proof obligation.\par
\textbf{Skeptic response:} This Skeptic response is a deterministic local fixture, not a live-model review. The exact six-state DP is a useful independent finite check, and the earlier counterexample correctly rejects only the simpler order-three guess. Agreement through n=48 does not imply a recurrence for all n. The report also needs to show the ordered block matrices and their characteristic relation, not just cite the candidate inherited from the roots-of-unity session.\par
\textbf{Assumptions reviewed:} The recurrence residual script uses exact Python integers.; No inference from finite ratios or finite recurrence checks is treated as proof.\par
\textbf{Proof obligations considered:} Show the ordered transfer-matrix polynomial identity symbolically.; Establish a rigorous asymptotic bound for the nontrivial modes.\par
\textbf{Still open:} Derive the order-12 recurrence from exact phase-ordered transfer-matrix products for every n.; Verify all initial/boundary indices algebraically rather than only by DP.; Prove the nontrivial root-of-unity terms are negligible compared with F\_(n+2).; Compute the exact three-step matrix product for each phase and show it yields the claimed recurrence.; State enough initial values and check the n=0 boundary exactly.; Prove the asymptotic contribution of the two nontrivial residue modes is o(F\_(n+2)).\par
\textbf{Solution evidence references:} 5e15b337-adf1-483c-8ba6-7cd2de5f500b\par
\textbf{Unresolved objections:} The general recurrence has not been derived from the matrices in this artifact.; The requested limiting density remains unproved.\par
\textbf{Resolved review notes:} None recorded\par
\textbf{Skeptic objection:} Finite recurrence residuals through n=48 are corroboration, not a proof for arbitrary n.\par
\textbf{Skeptic objection:} The transfer matrices are periodic but do not commute in general; the three-step product order must be stated and verified.\par
\textbf{Skeptic objection:} Even a proved recurrence and a\_12 value would not by itself establish the requested limiting ratio without control of characteristic roots and coefficients.\par
\textbf{Saved verification code:} incomplete\par
\textbf{Strategy finding:} Periodic transfer matrix · partial\par
\textbf{Strategy execution status:} ended\par
\textbf{Strategy failure detail:} Stop budget exhausted\par
\textbf{Claim:} Deterministic fixture result (not live-model output): a finite exact DP through n=60 is consistent with the candidate order-12 step-three recurrence; the general recurrence and limiting ratio remain unproved in this record.\par
\textbf{Source artifacts:} 5e15b337-adf1-483c-8ba6-7cd2de5f500b\par
\textbf{Objection:} Finite recurrence residuals through n=48 are corroboration, not a proof for arbitrary n.\par
\textbf{Objection:} The transfer matrices are periodic but do not commute in general; the three-step product order must be stated and verified.\par
\textbf{Objection:} Even a proved recurrence and a\_12 value would not by itself establish the requested limiting ratio without control of characteristic roots and coefficients.\par
\textbf{Open obligation:} Derive the order-12 recurrence from exact phase-ordered transfer-matrix products for every n.\par
\textbf{Open obligation:} Verify all initial/boundary indices algebraically rather than only by DP.\par
\textbf{Open obligation:} Prove the nontrivial root-of-unity terms are negligible compared with F\_(n+2).\par
\textbf{Open obligation:} Compute the exact three-step matrix product for each phase and show it yields the claimed recurrence.\par
\textbf{Open obligation:} State enough initial values and check the n=0 boundary exactly.\par
\textbf{Open obligation:} Prove the asymptotic contribution of the two nontrivial residue modes is o(F\_(n+2)).\par
\subsubsection*{Strategy assessments}
\begin{longtable}{>{\raggedright\arraybackslash}p{.16\linewidth}>{\raggedright\arraybackslash}p{.15\linewidth}>{\raggedright\arraybackslash}p{.40\linewidth}>{\raggedright\arraybackslash}p{.17\linewidth}}
\toprule Strategy & Assessment & Rationale & Evidence \\
\midrule\endhead
Periodic transfer matrix & Productive & The later attempt explicitly incorporates the earlier counterexample and adds an independent exact finite check while preserving the outstanding proof requirements. & 5e15b337-adf1-483c-8ba6-7cd2de5f500b \\
Roots-of-unity weighted Fibonacci polynomials & Not used & No session artifact is linked to this strategy. & No artifact references \\
Gap removal and Gaussian binomial sums & Not used & No session artifact is linked to this strategy. & No artifact references \\
\bottomrule
\end{longtable}
\paragraph{Artifact 1 · Unverified research note. Any Investigator Python or SageMath execution, including its source, status, output, and captured artifacts, is recorded below and is not proof by itself.}
\textbf{Claim:} Revisit the six-state transfer strategy using the earlier counterexample and later independent encoding: exact DP through n=60 agrees with the candidate order-12, step-three recurrence, but this finite check does not establish it for all n.\par
\textbf{Method:} Earlier deterministic fixture sessions recorded: Gap removal and Gaussian binomial sums: The gap-removal bijection gives P\_n(q)=sum\_k q\textasciicircum{}(k\textasciicircum{}2) [n-k+1 choose k]\_q; exact polynomial checks agree with the weighted recurrence through n=18.; Roots-of-unity weighted Fibonacci polynomials: Deterministic fixture result (not live-model output): the root-of-unity filter suggests a\_(n+12)-5a\_(n+9)+2a\_(n+6)+5a\_(n+3)+a\_n=0 and gives a\_12=128.; Periodic transfer matrix: A tentative three-step rule a\_(n+3)=a\_(n+2)+a\_n fits the first two transitions, but has not been justified. Session 1's explicit n=2 counterexample ruled out a\_(n+3)=a\_(n+2)+a\_n, so this pass does not revive that rule. Sessions 2 and 3 supplied a candidate order-12 relation and an independent q-binomial encoding. I implemented a direct six-state dynamic program and checked the candidate recurrence for every n from 0 through 48. This is a deterministic fixture narrative and computation, not live-model inference. The next proof step is to derive the relation symbolically from phase-ordered transfer matrices and handle the initial boundary separately.\par
\textbf{Assumptions:} The DP state stores the residue modulo 3 and whether the preceding integer was selected.; Indices are one-based and the empty subset contributes at n=0.; Prior session summaries are evidence records from deterministic local fixtures, not live research results.\par
\textbf{Evidence status:} Deterministic local fixture response; no live-model inference. A real local Python sandbox run checks exact integer counts through n=60; finite residuals do not prove the recurrence globally.\par
\textbf{Skeptic response:} This Skeptic response is a deterministic local fixture, not a live-model review. The exact six-state DP is a useful independent finite check, and the earlier counterexample correctly rejects only the simpler order-three guess. Agreement through n=48 does not imply a recurrence for all n. The report also needs to show the ordered block matrices and their characteristic relation, not just cite the candidate inherited from the roots-of-unity session.\par
\textbf{Objections:} Finite recurrence residuals through n=48 are corroboration, not a proof for arbitrary n.; The transfer matrices are periodic but do not commute in general; the three-step product order must be stated and verified.; Even a proved recurrence and a\_12 value would not by itself establish the requested limiting ratio without control of characteristic roots and coefficients.\par
\textbf{Decision:} stop\_budget\_exhausted\par
\textbf{Open obligations:} Derive the order-12 recurrence from exact phase-ordered transfer-matrix products for every n.; Verify all initial/boundary indices algebraically rather than only by DP.; Prove the nontrivial root-of-unity terms are negligible compared with F\_(n+2).; Compute the exact three-step matrix product for each phase and show it yields the claimed recurrence.; State enough initial values and check the n=0 boundary exactly.; Prove the asymptotic contribution of the two nontrivial residue modes is o(F\_(n+2)).\par
\textbf{Next question:} Not recorded\par
\textbf{Terminal reason:} stop\_budget\_exhausted\par
\subsection{Session \#5}
\begin{description}[style=nextline,font=\normalfont\bfseries,leftmargin=0pt,labelsep=0pt]
\item[Execution status] ended
\item[Outcome] Mixed · partial progress, unresolved
\item[Strategy] Gap removal and Gaussian binomial sums
\item[Started] 2026-10-23T21:31:06.111301+00:00
\end{description}
\textbf{Outcome summary:} Deterministic fixture result (not live-model output): the standard q-binomial shift corrects the previous exponent explanation, and exact polynomials agree through n=18; the global recurrence and limit remain unproved.\par
\subsection{Conclusion}
\#\#\# Historical session

No Conclusion was saved for this session. Its existing outcome summary and detailed evidence remain available below.

Structured argument steps were not recorded; saved narrative is available for inspection in Reasoning / proof and the research artifacts. Its completeness and correctness remain unassessed.\par
Outcome caveats:\par
\begin{itemize}
\item Agent-reported evidence is unverified; agreement is not independent verification.
\item The saved argument remains incomplete; no solution is claimed.
\item Finite exact checks through n=18 do not prove the all-n generating identity.
\item No complete recurrence or asymptotic proof is recorded.
\item Open obligations remain; inspect the saved artifacts before relying on this result.
\end{itemize}
\textbf{Verification:} Incomplete\par
Structured argument steps were not recorded; saved narrative is available for inspection in Reasoning / proof and the research artifacts. Its completeness and correctness remain unassessed.\par
\subsection{Outcome metrics}
\subsection{Operational metrics}
\textbf{Participating agents:} Investigator, Skeptic\par
\textbf{Model calls by status:} 2 dispatched · 2 completed · 0 rejected · 0 timed out · 0 interrupted / cancelled · 0 failed · 0 in progress · 0 not dispatched · 0 status unavailable\par
Tokens · input / output / cached / total: — (0/2 calls report token usage) / — (0/2 calls report token usage) / — (0/2 calls include cached-token metadata) / — (token usage incomplete for 2 call(s))\par
Wall-clock / agent call time: 0s wall · Investigator 0s across 1 call(s) · Skeptic 0s across 1 call(s)\par
\textbf{Python runs by status:} 1 completed\par
\textbf{SageMath runs by status:} 0 executions\par
\textbf{Artifacts and storage:} 1 finding(s) · 0 Python file(s) (0 B) · session record 56,344 B\par
Provider failures / retries: 0 provider failure(s) (0 rejected, 0 timed out, 0 other) · 0 interrupted / cancelled · 0 later attempt(s); no automatic retries\par
\textbf{Estimated monetary cost:} — (complete usage metadata or applicable price unavailable)\par
\subsection{Research-quality metrics}
\textbf{Strategies investigated:} Gap removal and Gaussian binomial sums\par
\textbf{Research cycles:} 1 completed / 1 started\par
\textbf{Outcome classification:} Mixed · partial progress, unresolved\par
\textbf{Verification status:} Incomplete\par
\textbf{Open proof obligations:} 7 · Prove the generating identity for all n and k under the stated convention.; Derive a general recurrence and establish the density limit. (+5 more)\par
\textbf{Evidence-backed conclusions:} 1 conclusion(s) · 1 cited finding(s)\par
\textbf{Strategy contribution:} not used: 2 · productive: 1\par
\textbf{Unresolved objections:} 2 · Finite exact checks through n=18 do not prove the all-n generating identity.; No complete recurrence or asymptotic proof is recorded.\par
\textbf{Claim scope:} not recorded\par
\textbf{Solution claim:} Deterministic fixture result (not live-model output): the standard q-binomial shift corrects the previous exponent explanation, and exact polynomials agree through n=18; the global recurrence and limit remain unproved.\par
\textbf{Argument:} The exponent correction reconciles the exact bijection with the standard Gaussian polynomial convention. It improves the saved reasoning, but the finite test is only corroboration and the main recurrence/limit proof is still missing.\par
\textbf{Skeptic response:} This Skeptic response is a deterministic local fixture, not a live-model review. The Investigator correctly distinguishes the formula from its earlier explanation: the standard Gaussian-binomial convention uses the offset k(k+1)/2 for subsets of \{1,...,m\}, so the total exponent is k(k-1)/2 + k(k+1)/2 = k\textasciicircum{}2. The exact polynomial comparison supports the identity for the tested finite range. The previous artifact's prose error is repaired here, but finite equality checks do not establish the identity for all n, the requested recurrence, or the density limit.\par
\textbf{Assumptions reviewed:} The q-binomial is partition-normalized and enumerates k-subsets with the baseline sum k(k+1)/2 removed.; The shift i\_j to y\_j is reversible and preserves the stated bounds.\par
\textbf{Proof obligations considered:} Prove the generating identity for all n and k under the stated convention.; Derive a general recurrence and establish the density limit.\par
\textbf{Still open:} Derive the required constant-coefficient recurrence globally from a justified root-of-unity or transfer-matrix identity.; Prove a\_n/F\_(n+2) tends to 1/3 with a bound on the nontrivial modes.; Rewrite the q-binomial derivation with the standard k(k+1)/2 normalization and prove it for every n.; Derive the constant-coefficient recurrence and sufficient initial values.; Prove the limiting ratio rather than extrapolating from finite coefficients.\par
\textbf{Solution evidence references:} 516c2288-3458-40b2-9e02-87a918c35a3a\par
\textbf{Unresolved objections:} Finite exact checks through n=18 do not prove the all-n generating identity.; No complete recurrence or asymptotic proof is recorded.\par
\textbf{Resolved review notes:} None recorded\par
\textbf{Skeptic objection:} The standard q-binomial identity should be stated with its precise subset-indexing convention; the previous artifact's explanatory factor was incorrect.\par
\textbf{Skeptic objection:} The polynomial equality is tested only for n\textless{}=18 and still needs a symbolic bijection/generating-function proof.\par
\textbf{Skeptic objection:} The recurrence and asymptotic density remain separate open goals.\par
\textbf{Saved verification code:} incomplete\par
\textbf{Strategy finding:} Gap removal and Gaussian binomial sums · partial\par
\textbf{Strategy execution status:} ended\par
\textbf{Strategy failure detail:} Stop budget exhausted\par
\textbf{Claim:} Deterministic fixture result (not live-model output): the standard q-binomial shift corrects the previous exponent explanation, and exact polynomials agree through n=18; the global recurrence and limit remain unproved.\par
\textbf{Source artifacts:} 516c2288-3458-40b2-9e02-87a918c35a3a\par
\textbf{Objection:} The standard q-binomial identity should be stated with its precise subset-indexing convention; the previous artifact's explanatory factor was incorrect.\par
\textbf{Objection:} The polynomial equality is tested only for n\textless{}=18 and still needs a symbolic bijection/generating-function proof.\par
\textbf{Objection:} The recurrence and asymptotic density remain separate open goals.\par
\textbf{Open obligation:} Derive the required constant-coefficient recurrence globally from a justified root-of-unity or transfer-matrix identity.\par
\textbf{Open obligation:} Prove a\_n/F\_(n+2) tends to 1/3 with a bound on the nontrivial modes.\par
\textbf{Open obligation:} Rewrite the q-binomial derivation with the standard k(k+1)/2 normalization and prove it for every n.\par
\textbf{Open obligation:} Derive the constant-coefficient recurrence and sufficient initial values.\par
\textbf{Open obligation:} Prove the limiting ratio rather than extrapolating from finite coefficients.\par
\subsubsection*{Strategy assessments}
\begin{longtable}{>{\raggedright\arraybackslash}p{.16\linewidth}>{\raggedright\arraybackslash}p{.15\linewidth}>{\raggedright\arraybackslash}p{.40\linewidth}>{\raggedright\arraybackslash}p{.17\linewidth}}
\toprule Strategy & Assessment & Rationale & Evidence \\
\midrule\endhead
Periodic transfer matrix & Not used & No session artifact is linked to this strategy. & No artifact references \\
Roots-of-unity weighted Fibonacci polynomials & Not used & No session artifact is linked to this strategy. & No artifact references \\
Gap removal and Gaussian binomial sums & Productive & This follow-up repairs a concrete exponent-convention error while preserving the global recurrence and limit as open obligations. & 516c2288-3458-40b2-9e02-87a918c35a3a \\
\bottomrule
\end{longtable}
\paragraph{Artifact 1 · Unverified research note. Any Investigator Python or SageMath execution, including its source, status, output, and captured artifacts, is recorded below and is not proof by itself.}
\textbf{Claim:} The earlier gap-removal formula P\_n(q)=sum\_k q\textasciicircum{}(k\textasciicircum{}2) [n-k+1 choose k]\_q is consistent with the standard Gaussian-binomial convention; the prior explanation of its exponent shift needs correction.\par
\textbf{Method:} Revisiting the saved gap-removal artifact 02ca743b-d66f-41c0-8692-34579ea5adfa: the position shift y\_j=i\_j-(j-1) gives sum(i\_j)=sum(y\_j)+k(k-1)/2. Under the standard convention, [m choose k]\_q sums q\textasciicircum{}(sum(y\_j)-k(k+1)/2) over 1\textless{}=y\_1\textless{}...\textless{}y\_k\textless{}=m, so the generating sum for q\textasciicircum{}(sum y\_j) contributes q\textasciicircum{}(k(k+1)/2), not q\textasciicircum{}(k(k-1)/2). Multiplying the two shifts gives q\textasciicircum{}(k\textasciicircum{}2), so the displayed formula survives while the earlier explanation is corrected. I then compare exact coefficient polynomials through n=18. This is deterministic local fixture research, not live-model inference.\par
\textbf{Assumptions:} The Gaussian polynomial uses the standard partition-normalized convention.; The positions y\_j range over 1 through n-k+1 and are strictly increasing.; The k=0 term is one and is handled separately in the finite check.\par
\textbf{Evidence status:} Deterministic fixture response; exact Python integer-polynomial comparison through n=18 is a finite check, not a proof of the limiting ratio.\par
\textbf{Skeptic response:} This Skeptic response is a deterministic local fixture, not a live-model review. The Investigator correctly distinguishes the formula from its earlier explanation: the standard Gaussian-binomial convention uses the offset k(k+1)/2 for subsets of \{1,...,m\}, so the total exponent is k(k-1)/2 + k(k+1)/2 = k\textasciicircum{}2. The exact polynomial comparison supports the identity for the tested finite range. The previous artifact's prose error is repaired here, but finite equality checks do not establish the identity for all n, the requested recurrence, or the density limit.\par
\textbf{Objections:} The standard q-binomial identity should be stated with its precise subset-indexing convention; the previous artifact's explanatory factor was incorrect.; The polynomial equality is tested only for n\textless{}=18 and still needs a symbolic bijection/generating-function proof.; The recurrence and asymptotic density remain separate open goals.\par
\textbf{Decision:} stop\_budget\_exhausted\par
\textbf{Open obligations:} Derive the required constant-coefficient recurrence globally from a justified root-of-unity or transfer-matrix identity.; Prove a\_n/F\_(n+2) tends to 1/3 with a bound on the nontrivial modes.; Rewrite the q-binomial derivation with the standard k(k+1)/2 normalization and prove it for every n.; Derive the constant-coefficient recurrence and sufficient initial values.; Prove the limiting ratio rather than extrapolating from finite coefficients.\par
\textbf{Next question:} Not recorded\par
\textbf{Terminal reason:} stop\_budget\_exhausted\par
\section{Agentic runs}
No agentic runs are recorded.\par
\end{document}
